\newpage
\renewcommand{\sectioncolor}{chap5}
\section{Konteneryzacja i środowisko uruchomieniowe}

W celu zapewnienia powtarzalności środowiska oraz uproszczenia procesu wdrażania, system został w pełni skonteneryzowany przy użyciu technologii \textbf{Docker} \cite{docker_docs}. Pozwala to na izolację poszczególnych mikroserwisów wraz z ich zależnościami w lekkich kontenerach, co eliminuje problem różnic między środowiskami programistycznymi a produkcyjnymi. Obrazy wynikowe są przechowywane i zarządzane w ramach rejestru \textit{GitHub Container Registry} \cite{ghcr_guide}.

\subsection{Konfiguracja wielokontenerowa (Docker Compose)}
Do zarządzania całym ekosystemem usług wykorzystano narzędzie \textbf{Docker Compose}. Pozwala ono na zdefiniowanie wszystkich komponentów systemu (mikroserwisy, bazy danych, brokerzy komunikatów) w jednym pliku konfiguracyjnym \texttt{docker-compose.yml} oraz ich jednoczesne uruchomienie przy pomocy jednej komendy.

\begin{lstlisting}[language=bash, caption={Uruchomienie całego stosu technologicznego systemu}]
# Budowa obrazów i uruchomienie kontenerów w tle
docker-compose up --build -d
\end{lstlisting}

\subsection{Charakterystyka zdefiniowanych usług}
Zgodnie z architekturą systemu, w pliku konfiguracyjnym zdefiniowano następujące grupy usług:

\begin{itemize}
    \item \textbf{Infrastruktura danych (PostgreSQL):} Wykorzystano obraz \textit{postgres:15-alpine} \cite{postgresql_docs}. Kontener został skonfigurowany tak, aby przy starcie automatycznie tworzył trzy odrębne bazy danych (\texttt{identitydb}, \texttt{reservationdb}, \texttt{notificationdb}) za pomocą skryptu inicjalizacyjnego. Trwałość danych zapewniono poprzez wolumeny (\textit{Volumes}).
    \item \textbf{Broker komunikatów (RabbitMQ):} Zastosowano obraz z panelem zarządzania (\textit{management-alpine}) \cite{rabbitmq_docs}, co pozwala na monitorowanie kolejek zdarzeń w czasie rzeczywistym na porcie 15672.
    \item \textbf{Warstwa mikroserwisów (.NET):} Serwisy \textit{Identity}, \textit{Reservation} oraz \textit{Notification} są budowane z lokalnych plików \texttt{Dockerfile} w oparciu o SDK .NET 8 \cite{dotnet_microsoft}. 
    \item \textbf{API Gateway (Ocelot):} Wykorzystano bibliotekę Ocelot \cite{ocelot_docs}, która pełni rolę centralnego punktu wejścia, mapując zewnętrzny port 5000 na wewnętrzną infrastrukturę kontenerów.
    \item \textbf{Narzędzia pomocnicze (MailHog):} Kontener służący do przechwytywania i testowania wiadomości e-mail w środowisku deweloperskim, co pozwala na bezpieczną weryfikację systemu powiadomień bez wysyłki realnych wiadomości.
\end{itemize}

[Image of Docker Compose architecture for microservices showing PostgreSQL, RabbitMQ, and .NET containers connected via a bridge network]

\subsection{Zarządzanie zależnościami i Healthchecks}
W konfiguracji zastosowano mechanizm \textbf{Healthcheck}, który monitoruje stan gotowości usług krytycznych. Mikroserwisy posiadają instrukcję \texttt{depends\_on} z warunkiem \texttt{service\_healthy}, co gwarantuje, że aplikacja .NET spróbuje nawiązać połączenie z bazą dopiero wtedy, gdy silnik bazy danych przejdzie pomyślnie test gotowości.

\begin{lstlisting}[language=yaml, caption={Fragment konfiguracji Docker Compose obsługujący zależności}]
reservation_api:
  build:
    context: ./ReservationService
  depends_on:
    postgres_db:
      condition: service_healthy
    rabbitmq:
      condition: service_healthy
\end{lstlisting}

\subsection{Izolacja sieciowa}
Wszystkie kontenery działają w ramach dedykowanej sieci wirtualnej typu \textit{bridge}. Pozwala to na komunikację między serwisami przy użyciu nazw ich usług (np. \texttt{Host=postgres\_db}) zamiast zmiennych adresów IP. Takie rozwiązanie zapewnia bezpieczną izolację — bazy danych nie są wystawione na świat zewnętrzny, a dostęp do nich mają wyłącznie autoryzowane mikroserwisy wewnątrz sieci Docker.

\subsection{Automatyzacja zarządzania systemem (Skrypty Manage)}
W celu maksymalnego uproszczenia pracy z systemem, przygotowano dedykowane skrypty zarządcze: \texttt{manage.bat} (dla systemu Windows) oraz \texttt{manage.sh} (dla systemów Unix). Stanowią one warstwę abstrakcji nad surowymi komendami silnika Docker oraz narzędzia Vite \cite{vite_guide}.

\begin{lstlisting}[language=bash, caption={Uruchomienie kompletnego środowiska za pomocą skryptu zarządczego}]
# Uruchomienie backendu w kontenerach oraz frontendu React
manage.bat all
\end{lstlisting}

Wykorzystanie powyższego skryptu wykonuje sekwencyjnie następujące zadania:
\begin{itemize}
    \item \textbf{Weryfikacja środowiska:} Sprawdzenie dostępności niezbędnych narzędzi (Docker, Node.js).
    \item \textbf{Orkiestracja kontenerów:} Wywołanie \texttt{docker-compose} w celu zbudowania i uruchomienia infrastruktury.
    \item \textbf{Uruchomienie warstwy prezentacji:} Automatyczny start serwera deweloperskiego Vite \cite{vite_guide} dla aplikacji frontendowej.
\end{itemize}

Dzięki tej automatyzacji proces wdrożenia i testowania (w tym skanowanie podatności obrazów narzędziem Trivy \cite{trivy_security}) został zintegrowany z procesem ciągłej integracji w GitHub Actions \cite{github_actions_docs}.